\section{研究问题}

\begin{frame}{为什么需要“心理学版图灵测试”？}
  \begin{block}{核心问题}
    当前 LLM 评测多关注规范正确性、任务准确率或通用能力表现。
    但社会心理学中的许多经典现象并不是“标准答案”问题，而是人类在有限理性与社会情境下形成的稳定行为模式。
  \end{block}

  \vspace{0.6em}
  \begin{columns}[T,totalwidth=\textwidth]
    \begin{column}{0.48\textwidth}
      \textbf{传统基准更容易回答}
      \begin{itemize}
        \item 模型是否答对？
        \item 模型是否遵循规范推理？
        \item 模型是否在任务上更强？
      \end{itemize}
    \end{column}
    \begin{column}{0.48\textwidth}
      \textbf{本工作关注}
      \begin{itemize}
        \item 模型是否做得像人类？
        \item 模型是否对社会线索敏感？
        \item 模型是否具有相似的心理价值权重？
      \end{itemize}
    \end{column}
  \end{columns}
\end{frame}

\begin{frame}{目标：检验 LLM 的社会理性}
  \begin{block}{社会理性}
    本文将社会理性理解为：在具有心理适应性机制的社会决策情境中，决策者对亲缘关系、群体规模、收益变化、身份线索与自我相关前提等因素表现出的系统性敏感性。
  \end{block}

  \vspace{0.7em}
  \begin{itemize}
    \item 不是问 LLM 是否给出“最优”答案。
    \item 而是问 LLM 是否复现人类在经典范式中的\metric{类人决策模式}。
    \item 评估对象从单题准确率转向\metric{行为分布}与\metric{条件变化下的趋势结构}。
  \end{itemize}
\end{frame}

\begin{frame}{本文的基准尝试弥补什么空缺？}
  \begin{columns}[T,totalwidth=\textwidth]
    \begin{column}{0.50\textwidth}
      \begin{block}{现有评测的不足}
        \begin{itemize}
          \item 较少覆盖人类非规范性判断。
          \item 较少显式比较情境敏感性。
          \item 较少把社会认知偏差转化为可复现指标。
        \end{itemize}
      \end{block}
    \end{column}
    \begin{column}{0.46\textwidth}
      \begin{block}{本框架的定位}
        \begin{itemize}
          \item 复现经典社会心理学范式。
          \item 将人类基线写入结构化行为测试。
          \item 自动生成测试特定的行为指标与可视化。
        \end{itemize}
      \end{block}
    \end{column}
  \end{columns}

  \vspace{0.8em}
  \centering
  \Large\metric{从“能不能答对”转向“是不是像人类那样判断”。}
\end{frame}
